\documentclass{article}

\usepackage{booktabs}
\usepackage{tabularx}
\usepackage{hyperref}

\hypersetup{
    colorlinks=true,       % false: boxed links; true: colored links
    linkcolor=red,          % color of internal links (change box color with linkbordercolor)
    citecolor=green,        % color of links to bibliography
    filecolor=magenta,      % color of file links
    urlcolor=cyan           % color of external links
}

\title{Hazard Analysis\\\progname}

\author{\authname}

\date{}

\input{../Comments.text}
\input{../Common.text}

\begin{document}

\maketitle
\thispagestyle{empty}

~\newpage

\pagenumbering{roman}

\begin{table}[hp]
\caption{Revision History} \label{TblRevisionHistory}
\begin{tabularx}{\textwidth}{llX}
\toprule
\textbf{Date} & \textbf{Developer(s)} & \textbf{Change}\\
\midrule
Oct 16, 2026 & Farid Bastoros & Initial revision of 1 Introduction.\\
Date1 & Name(s) & Description of changes\\
Date2 & Name(s) & Description of changes\\
... & ... & ...\\
\bottomrule
\end{tabularx}
\end{table}

~\newpage

\tableofcontents

\listoffigures

\listoftables

~\newpage

\pagenumbering{arabic}

\wss{You are free to modify this template.}

\section{Introduction} \label{SecIntro}

% \wss{You can include your definition of what a hazard is here.}

\progname{} is a mobile application in which users earn virtual currency from
the physical activity their phone records, invest that currency in a simulated
market that follows real-world prices, learn the financial and fitness concepts
behind their decisions, and share chosen progress with friends. Although
\progname{} controls no physical hardware and never moves real money, it
handles health data, presents financial information that users may act on, and
rewards physical exertion. Failures in any of these areas can harm users, so a
systematic analysis of what can go wrong is needed before the software is
designed.

This document identifies the hazards of \progname{}, analyzes how each could
arise, and derives requirements that eliminate or reduce them. It uses the
following definitions:

\begin{itemize}
\item A \textbf{hazard} is a state or set of conditions of \progname{} that,
  together with a set of worst-case conditions in its environment, will lead to
  a loss~\cite{Leveson2011}. For example, a
  shared update that reveals a user's weight is a hazard: combined with a
  friend who sees it, it leads to a loss of privacy.
\item A \textbf{loss} is harm to a stakeholder that the team wants to prevent.
  The losses considered in this document are described in
  Section~\ref{SecScopePurpose}.
\item A \textbf{failure mode} is a way in which a component of \progname{} can
  fail to perform its intended function~\cite{IEC60812}. A component fails not
  only when its output is missing, but also when its output is wrong, late, or
  duplicated; this analysis considers all four.
\item A \textbf{safety or security requirement} is a requirement that
  eliminates a hazard, reduces its likelihood, or reduces its effects.
\end{itemize}

The rest of this document is organized as follows.
Section~\ref{SecScopePurpose} states the scope and purpose of the analysis and
the losses it aims to prevent. Section~\ref{SecBoundaries} defines the system
boundary and divides \progname{} into components.
Section~\ref{SecAssumptions} lists the critical assumptions.
Section~\ref{SecFMEA} presents the Failure Mode and Effect Analysis (FMEA),
which is the core of this document. Section~\ref{SecSafetyReqs} collects the
resulting safety and security requirements, and Section~\ref{SecRoadmap}
states when each will be implemented.

\section{Scope and Purpose of Hazard Analysis} \label{SecScopePurpose}

\wss{You should say what \textbf{loss} could be incurred because of the
hazards.}

\section{System Boundaries and Components} \label{SecBoundaries}

\wss{Dividing the system into components will help you brainstorm the hazards.
You shouldn't do a full design of the components, just get a feel for the major
ones.  For projects that involve hardware, the components will typically include
each individual piece of hardware.  If your software will have a database, or an
important library, these are also potential components.}

\section{Critical Assumptions} \label{SecAssumptions}

\wss{These assumptions that are made about the software or system.  You should
minimize the number of assumptions that remove potential hazards.  For instance,
you could assume a part will never fail, but it is generally better to include
this potential failure mode.}

\section{Failure Mode and Effect Analysis} \label{SecFMEA}

\wss{Include your FMEA table here. This is the most important part of this document.}
\wss{The safety requirements in the table do not have to have the prefix SR.
The most important thing is to show traceability to your SRS. You might trace to
requirements you have already written, or you might need to add new
requirements.}
\wss{If no safety requirement can be devised, other mitigation strategies can be
entered in the table, including strategies involving providing additional
documentation, and/or test cases.}

\section{Safety and Security Requirements} \label{SecSafetyReqs}

\wss{Newly discovered requirements.  These should also be added to the SRS.  (A
rationale design process how and why to fake it.)}

\section{Roadmap} \label{SecRoadmap}

\wss{Which safety requirements will be implemented as part of the capstone timeline?
Which requirements will be implemented in the future?}

\newpage{}

\section{Appendix --- Generative AI Use Disclosure}

\wss{ChatGPT (version 5) was used to brainstorm ideas for the template for this disclose section.}

\subsection{AI Tools Used}

\wss{Give the name of the AI tool(s) used and the version.}

\subsection{How and Where Used}

\wss{Explain what the tool(s) were used for.  Examples include: brainstorming,
explaining a technical concept, reviewing the document, generating or modifying
text, generating or modifying diagrams, identifying errors or omissions or
inconsistencies.}

\wss{For each of the uses, identify which parts of the document were affected.}

\wss{You do not need to report every interaction, but you should disclose the 
uses that materially contributed to your work}

\subsection{Verification of AI-Generated Content}

\wss{\begin{itemize}
    \item Describe how the team checked AI-generated or AI-assisted material
    before including it in the document.
    \item In particular, verify factual claims, technical assertions,
    requirements, calculations, references, and other information for which an
    AI system may produce incorrect or fabricated results.
    \item \textbf{The team is responsible for the accuracy, completeness,
consistency, and appropriateness of the submitted document, regardless of
whether material was generated by a human or an AI system.}
\end{itemize}}

\newpage{}

\section{Appendix --- Reflection}

\wss{Not required for CAS 741}

\input{../Reflection.text}

\begin{enumerate}
    \item What went well while writing this deliverable? 
    \item What pain points did you experience during this deliverable, and how
    did you resolve them?
    \item Which of your listed risks had your team thought of before this
    deliverable, and which did you think of while doing this deliverable? For
    the latter ones (ones you thought of while doing the Hazard Analysis), how
    did they come about?
    \item Other than the risk of physical harm (some projects may not have any
    appreciable risks of this form), list at least 2 other types of risk in
    software products. Why are they important to consider?
\end{enumerate}

\newpage{}

\addcontentsline{toc}{section}{References}
\begin{thebibliography}{99}

\bibitem{IEC60812}
International Electrotechnical Commission, \emph{IEC 60812:2018 Failure Modes
and Effects Analysis (FMEA and FMECA)}, 3rd ed. Geneva, Switzerland: IEC,
2018.

\bibitem{Leveson2011}
N.~G. Leveson, \emph{Engineering a Safer World: Systems Thinking Applied to
Safety}. Cambridge, MA, USA: MIT Press, 2011.


\end{thebibliography}

\end{document}